\section{Introduction}\label{sec:intro}

When a frontier model ships, its provider publishes a model card carrying benchmark results, and independent evaluators run standardised suites over the same models \citep{liang2022helm,epoch2026hub}. This paper asks whether benchmark availability can identify whether the results a provider leaves out stand below the ones it prints, the older economic question. Disclosure unravels only when the sender's evidence is commonly known \citep{grossman1981,milgrom1981}. \citet{dye1985} shows what breaks that: a receiver who cannot tell a sender holding nothing from one holding something bad pools the two, and withholding survives. Any study of what was left out inherits that pooling, since the object it needs is the endowment, the results a sender held when it wrote, which the work below does not observe.

Benchmark documentation is one of the few settings where that endowment can be bounded from outside: absent prepublication sharing, a benchmark published after a model shipped cannot have been in its endowment, whatever a provider ran privately, so dating the instruments separates absence that could have been a disclosure from absence that could not.

Our first contribution is a panel built around that bound. From a pinned snapshot of an independent benchmarking hub covering 819 model-versions, we collapse scaffold variants to the release a provider ships, giving 353 releases across 46 organisations, 74 benchmarks and 3,082 scored pairs. Every benchmark carries a verified release date, splitting the panel into 2,355 eligible pairs and 476 that could not have been disclosed.

Our second contribution is that bounding the endowment is not enough, since judging withholding needs a measure of how results stood against peers. The eventual comparison is reported versus omitted benchmarks among those available, and it needs the two sides of availability balanced in standing. They are not. Benchmarks that already existed when a model shipped outscore benchmarks built afterwards by 12.23 percentile points on average, positive in 78.1\% of the 146 releases holding both kinds of cell, and nothing in that calculation knows which benchmarks a provider reported. The most favourable measure of standing we can construct still leaves 9.14 points. The contamination is structural, because the date identifying the endowment also shifts the measure of standing used to price silence \citep{heckman1979}. Read as economics: instrument vintages are a free verification technology for the endowment that \citet{dye1985} shows a receiver cannot otherwise see; the verifying variable is not excludable from the receiver's valuation of silence, so what it verifies it contaminates.

Our third contribution takes the question outside our own panel. Across 14 releases
of HELM Lite, the pool grows from 30 models to 91 while 24 models keep bit-identical scores, and every one of their published headline values moves, five pairs reversing order. A companion leaderboard from the same team, behind an absolute headline, shows none of this.

Our fourth contribution codes the record. Hand-coding the 100 releases with a locatable release-time document, blind-replicated on a fifth at $\kappa = 0.95$, yields the estimator such an audit invites: omitted benchmarks stand $+11.4$ points above those a provider could not have reported, excluding zero yet short of its matched label-free counterpart by only the wedge composition implies: it cannot distinguish concealment from the availability artifact. What survives is composition: reported benchmarks stand 8.7 points above omitted eligible ones within release ($p = 0.0001$), a gap benchmark fixed effects do not remove.

\subsection{Related literature}\label{sec:related}

That reporting practice in model documentation is uneven, incomplete and shaped by the sender is established, and we concede it \citep{staufer2025auditcards,kuehnert2026disclosure}. Three papers sit closest. \citet{zhang2026positional} audits four flagship releases inside a long-context evaluation study and finds a benchmark family absent from every main result table, an omission count without the dates that would tell a passed-over benchmark from one the release predates. \citet{ghosh2026evalcards} builds the reporting format in which an omission would become visible, and asks what a reader could see rather than what was left out. \citet{singh2025leaderboard} documents selection inside a live leaderboard, providers testing private variants and keeping the best, alongside evaluator policies that move the ranking; what it cannot separate, and we make separable at the cell level, is an omission from a result that could not have existed. \citet{bean2025measuring} reviews 445 benchmarks for construct validity, whereas we take the instruments as given and ask which results a provider prints. None of the four conditions on availability, the single thing we add: without a verified benchmark date, an omission and a nonexistence are the same observation. The economics runs from unravelling under a known endowment \citep{grossman1981,milgrom1981} to its failure under endowment uncertainty \citep{dye1985} and under costly disclosure \citep{verrecchia1983}, frictions this design does not separate. The report-card literature supplies the cautionary cases \citep{dranove2010,hotz2013,jin2003}, including withholding by high-quality senders \citep{bederson2018}. Concurrent theory models strategic evaluation and persuasion directly \citep{truong2026strategic,dai2026persuasion}; that work is theory; this paper is measurement.
